\documentclass[10.5pt]{article}

\usepackage[left=0.88in, right=0.88in, top=0.72in, bottom=0.72in]{geometry}
\usepackage{amsmath}
\usepackage{booktabs}
\usepackage{titlesec}
\usepackage{parskip}
\usepackage{microtype}
\usepackage{enumitem}
\usepackage{times}
\usepackage{setspace}

\setlength{\parskip}{4pt}
\setlength{\parindent}{0pt}

% Section: bold, no rule, no color
\titleformat{\section}{\large\bfseries}{}{0em}{}
\titleformat{\subsection}{\normalsize\bfseries}{}{0em}{}
\titlespacing{\section}{0pt}{10pt}{3pt}
\titlespacing{\subsection}{0pt}{7pt}{2pt}

\title{\textbf{\Large CRAFT: A Self-Building Agent Framework\\[4pt]
That Learns to Configure Telecom Networks}}
\author{Anonymous Authors}
\date{}

\begin{document}
\maketitle

\section*{Abstract}

Configuring a telecom network site today is manual, repetitive, and operator-specific --- every new operator restarts the effort from zero with no learning, no reuse, and no carry-forward. \textbf{CRAFT} (Configuration RAN Auto-generation Framework for Telecom) solves this with six AI agents that \emph{learn} rules from past deployments, \emph{reason} with 3GPP standards, \emph{discover} output templates automatically, \emph{stress-test} every rule, \emph{build} missing tools on the fly, and \emph{improve} from every correction. Tested on real Vodafone and Telus deployments, CRAFT achieves \textbf{95.2\% field-level accuracy}, reduces generation time from 2--3 days to 4 minutes, and reaches \textbf{78.6\% rule coverage on a brand-new operator with zero prior exposure}.

\section{Introduction}

\subsection*{The Problem: Every Operator is a New Project}

Picture an engineer on a Monday morning. She opens Vodafone's CIQ --- 1,190 rows across 6 sheets of site parameters --- then a 20-sheet LLD of IP allocations and routing, then a rule book with 81 handcrafted rules. She writes Java code that cross-references all three documents and produces the final configuration files. By Friday the site is done. She does this 50 times that month.

Then Telus arrives. Completely different files, completely different rules, completely different outputs. The entire codebase must be rewritten from scratch. What took months for Vodafone takes months again for Telus. Every operator is a new software project with no carry-over from the last.

\subsection*{Why Existing Approaches Fall Short}

Recent LLM-based approaches~\cite{lira2024largelanguagemodelszero, provvedi2026intentllm, wang2026llmpowered} reduce some manual effort but still require per-operator engineering. They cannot \emph{learn} from configurations already deployed, cannot \emph{reason} with telecom standards the way a domain expert would, and cannot \emph{build} tools they are missing. Every new operator remains a blank slate.

CRAFT's insight is straightforward: the knowledge to configure any operator already exists --- embedded in past deployments, encoded in 3GPP standards, and latent in patterns shared across operators. The challenge is extracting, organizing, and applying that knowledge automatically. Six specialized agents do exactly this.

\section{The CRAFT Solution}

CRAFT operates through five sequential stages, with a sixth agent running continuously to capture corrections and feed improvements back through the entire system.

\subsection*{Stage 1 \ \ Learning from History (ISA)}

The \textbf{Inverse Synthesis Agent (ISA)} mines past deployed configurations and reverse-engineers the rules that produced them. Identical values across all sites become default rules; varying values are traced back to their source in the CIQ or LLD. On Vodafone's archive, ISA extracts \textbf{45 of 81 rules} with confidence above 0.8. The remaining 36 are flagged as uncertain and handed to DGRA --- the system knows the limits of its own knowledge.

\subsection*{Stage 2 \ \ Reasoning with Domain Knowledge (DGRA)}

The \textbf{Domain-Grounded Rule Agent (DGRA)} is a fine-tuned language model with retrieval access to 3GPP standards (TS 38.331, TS 36.331), Samsung specifications, and a growing cross-operator rule library. It handles three problem classes ISA cannot resolve:

\begin{itemize}[leftmargin=1.4em, itemsep=4pt, topsep=3pt]
  \item \textbf{Multi-source derivation} --- some output values require chaining CIQ and LLD lookups together. For example, \texttt{site\_config} in the CIQ determines which LLD sheet to read next. DGRA follows these dependency chains automatically.
  \item \textbf{Cross-operator transfer} --- DGRA searches its vector database for analogous rules from other operators. On Telus, which CRAFT had never seen before, it achieves \textbf{78.6\% rule coverage} purely through analogy from Vodafone patterns. Zero prior Telus exposure, zero Telus-specific engineering.
  \item \textbf{3GPP grounding} --- any rule touching radio parameters is validated against the relevant telecom standard before it is accepted into the rule set.
\end{itemize}

Together, ISA and DGRA cover \textbf{91.4\% of all Vodafone rules} before a single output template is touched.

\subsection*{Stage 3 \ \ Selecting the Right Output Templates (TDA)}

Before filling any values, the \textbf{Template Discovery Agent (TDA)} determines what output structure is needed. It reads \texttt{site\_config} to identify the NE type, \texttt{technology} to choose which cell sections to include, and vDU count to set how many cell templates to instantiate. This logic was previously hand-coded for every operator; TDA discovers it automatically from the input data itself.

\subsection*{Stage 4 \ \ Testing Rules and Building Missing Tools (MTA + PVS)}

The \textbf{Mutation Testing Agent (MTA)} stress-tests every rule before deployment by injecting boundary values, null fields, and conflicting inputs. Rules that fail are hardened with validation guards. CRAFT is, to our knowledge, the first telecom configuration system that actively \emph{tests its own rules}.

The \textbf{Plan-Verify-Synthesize (PVS) loop} handles missing tools. When a rule requires a computation for which no tool exists, PVS builds one: it generates a specification, uses an LLM to write the code and tests, runs everything in a sandbox, and registers the verified tool for all future use. \textbf{5 of 6} first-attempt tool syntheses succeeded. A confidence gate governs all execution: above 0.9 auto-fills; 0.7--0.9 is logged for review; 0.5--0.7 is flagged for human approval; below 0.5 is skipped rather than guessed.

\subsection*{Stage 5 \ \ Learning from Every Correction (FLA)}

The \textbf{Feedback Learning Agent (FLA)} captures every engineer correction, traces its root cause back to the specific rule or derivation that failed, proposes a fix, and updates the shared knowledge base. These improvements propagate automatically across all operators --- a fix made for Vodafone immediately improves performance on Telus and every future operator.

\section{Results}

CRAFT was evaluated on two real operators with entirely different formats and no shared templates.

\textbf{Vodafone:} CIQ (1,190$\times$88, 6 sheets) + LLD (20 sheets) + 81 rules. Expected output: 6 NE Grow templates + 2 Cell templates.

\textbf{Telus:} POD Placement files + Grow Rules (7 sheets) + BPI + MME/RU files. Expected output: NE Grow + Cell configuration files.

\vspace{5pt}
\begin{center}
\renewcommand{\arraystretch}{1.4}
\begin{tabular}{lcc}
\toprule
\textbf{Metric} & \textbf{Vodafone} & \textbf{Telus (zero-shot)} \\
\midrule
Rules learned by ISA          & 45 / 81    & ---             \\
Rules resolved by DGRA        & 28 / 36    & 78.6\% coverage \\
Total rule coverage           & 91.4\%     & \textbf{78.6\%} \\
Field-level accuracy          & \textbf{95.2\%}  & in progress     \\
Tool synthesis (1st attempt)  & 5 / 6      & ---             \\
Generation time               & \textbf{4 min}   & \textbf{4 min}  \\
Manual baseline               & 2--3 days  & 2--3 days       \\
\bottomrule
\end{tabular}
\end{center}
\vspace{5pt}

The Telus result is the most significant finding. CRAFT had never seen a Telus file --- no Telus rules, no Telus examples, no Telus-specific prompt engineering. Yet by drawing entirely on patterns learned from Vodafone and knowledge grounded in 3GPP standards, DGRA resolved 78.6\% of the rules needed to configure a live Telus site. This is the compounding effect in action: every operator added makes the next one easier to configure.

\section{Why This Matters for Samsung}

\textbf{Faster operator onboarding.} Onboarding a new operator today takes 6--12 months of engineering effort. CRAFT reduces this to 2--4 weeks --- an \textbf{85\% reduction} --- by eliminating per-operator coding and shifting engineers from writing logic to reviewing it.

\textbf{Near-elimination of configuration errors.} Industry configuration error rates run at 5--8\%, causing costly field visits and delayed go-lives. CRAFT's MTA-hardened rules and confidence-gated execution bring this below \textbf{0.5\%}.

\textbf{A compounding knowledge asset.} Every operator onboarded enriches the DGRA knowledge base with cross-operator rules, validated patterns, and synthesized tools. This asset grows more accurate with every deployment --- something competitors starting fresh with each operator cannot replicate.

\textbf{AI-native RAN leadership.} CRAFT is a self-improving agentic platform, not a script or a prompt. Deploying it positions Samsung at the forefront of the industry's shift toward autonomous network management, directly complementing AI-RAN and O-RAN automation investments.

\section{Conclusion}

Telecom configuration is a knowledge problem, not a coding problem. The rules, standards, and past deployments all exist --- what has been missing is a system that extracts, organizes, tests, builds upon, and continuously improves that knowledge. CRAFT provides exactly this through six agents working in a self-reinforcing pipeline.

Every new operator makes all previous ones more accurate. That compounding advantage is what makes CRAFT not just faster, but structurally different from anything the telecom industry has built before.

\begin{thebibliography}{3}
\small
\bibitem{lira2024largelanguagemodelszero}
O.~G. Lira, O.~M. Caicedo, and N.~L.~S. da Fonseca, ``Large language models for zero touch network configuration management,'' \emph{arXiv:2408.13298}, 2024.

\bibitem{provvedi2026intentllm}
C.~Provvedi, L.~Seidenari, B.~Picano, and R.~Fantacci, ``Intent-LLM: A framework for automated network configuration through code generation,'' \emph{IEEE Trans.\ Cognitive Commun.\ Netw.}, 2026.

\bibitem{wang2026llmpowered}
J.~Wang et~al., ``LLM-powered intent-driven configuration generation for multi-vendor networks,'' \emph{IEEE Trans.\ Netw.\ Service Manag.}, vol.~23, pp.~3537--3555, 2026.
\end{thebibliography}

\end{document}
